\section{Available Inventory Identity: Eight-Scenario Verification}
\label{app:avail-identity}

The available-inventory identity (P20, \ref{inv:P20})
\[
  \mathrm{avail} \;=\; \mathrm{own} - \mathrm{onloan} + \mathrm{borr}
\]
is verified across eight scenarios. Each scenario exercises a distinct combination of the six stored coordinates: own, onl, borr, c\_p ($\mathrm{coll\_post}$), c\_r ($\mathrm{coll\_recv}$), c\_rh ($\mathrm{coll\_rehyp}$). Available inventory is never stored. It is computed on read from three of those coordinates: own, onl, borr. The identity cannot be violated by construction.

Each row below states the stored coordinates for one entity in one scenario. The final column recomputes available inventory from those coordinates and checks the result.

\begin{center}\footnotesize
\setlength{\tabcolsep}{4pt}
\begin{longtable}{l r r r r r r l}
\toprule
\textbf{Entity (unit)} & own & onl & borr & c\_p & c\_r & c\_rh & $\mathrm{avail}=\mathrm{own}-\mathrm{onl}+\mathrm{borr}$ \\
\midrule
\endhead
\multicolumn{8}{@{}p{\dimexpr\textwidth-2\tabcolsep\relax}@{}}{\textit{A. Alice buys 1{,}000 AAPL outright.}} \\
Alice (AAPL) & 1000 & 0 & 0 & 0 & 0 & 0 & $1000-0+0=1000$ \checkmark \\
Virtual & $-1000$ & --- & --- & --- & --- & --- & (issuance counterpart) \\
\midrule
\multicolumn{8}{@{}p{\dimexpr\textwidth-2\tabcolsep\relax}@{}}{\textit{B. Alice owns 1{,}000 VOD, lends 400 to Bob.}} \\
Alice (VOD) & 1000 & 400 & 0 & 0 & 0 & 0 & $1000-400+0=600$ \checkmark \\
Bob (VOD) & 0 & 0 & 400 & 0 & 0 & 0 & $0-0+400=400$ \checkmark \\
\midrule
\multicolumn{8}{@{}p{\dimexpr\textwidth-2\tabcolsep\relax}@{}}{\textit{C. Bob lends 1{,}000 NVDA to Alice, Alice sells 500.}} \\
Bob (NVDA) & 1000 & 1000 & 0 & 0 & 0 & 0 & $1000-1000+0=0$ \checkmark \\
Alice (NVDA) & $-500$ & 0 & 1000 & 0 & 0 & 0 & ${-}500-0+1000=500$ \checkmark \\
\midrule
\multicolumn{8}{@{}p{\dimexpr\textwidth-2\tabcolsep\relax}@{}}{\textit{D. Swap desk holds 5Y EUR IRS (non-lendable).}} \\
Desk (IRS) & $+1$ & 0 & 0 & 0 & 0 & 0 & N/A (non-lendable; $\mathrm{own}=1$) \\
\midrule
\multicolumn{8}{@{}p{\dimexpr\textwidth-2\tabcolsep\relax}@{}}{\textit{E. Bob lends 1{,}000 NVDA, Alice sells all 1{,}000.}} \\
Bob (NVDA) & 1000 & 1000 & 0 & 0 & 0 & 0 & $1000-1000+0=0$ \checkmark \\
Alice (NVDA) & $-1000$ & 0 & 1000 & 0 & 0 & 0 & ${-}1000-0+1000=0$ \checkmark \\
\midrule
\multicolumn{8}{@{}p{\dimexpr\textwidth-2\tabcolsep\relax}@{}}{\textit{F. Alice holds 1{,}000 AAPL, lends 400 to Bob, Bob sells 200.}} \\
Alice (AAPL) & 1000 & 400 & 0 & 0 & 0 & 0 & $1000-400+0=600$ \checkmark \\
Bob (AAPL) & $-200$ & 0 & 400 & 0 & 0 & 0 & ${-}200-0+400=200$ \checkmark \\
\midrule
\multicolumn{8}{@{}p{\dimexpr\textwidth-2\tabcolsep\relax}@{}}{\textit{G. On-lending chain: Alice lends 600 MSFT to Bob, Bob re-lends 300 to Charlie.}} \\
Alice (MSFT) & 1000 & 600 & 0 & 0 & 0 & 0 & $1000-600+0=400$ \checkmark \\
Bob (MSFT) & 0 & 300 & 600 & 0 & 0 & 0 & $0-300+600=300$ \checkmark \\
Charlie (MSFT) & 0 & 0 & 300 & 0 & 0 & 0 & $0-0+300=300$ \checkmark \\
\midrule
\multicolumn{8}{@{}p{\dimexpr\textwidth-2\tabcolsep\relax}@{}}{\textit{H. Alice lends 400 AAPL to Bob under security interest; Bob posts 200 MSFT as pledge.}} \\
Bob (MSFT) & 0 & 0 & 0 & 200 & 0 & 0 & $0-0+0=0$ \checkmark \\
Alice (MSFT) & 0 & 0 & 0 & 0 & 200 & 0 & $0-0+0=0$ \checkmark \\
\bottomrule
\end{longtable}
\end{center}

Scenario G is the case that separates the correct projection from its naive alternative.

\begin{remark}[On-lending defeats the naive identity]
For intermediary re-lender Bob (G), $\mathrm{own}=0 \neq \mathrm{avail}+\mathrm{onl}=300+300$: the identity $\mathrm{own}=\mathrm{avail}+\mathrm{onloan}$ fails. The correct projection $\mathrm{avail}=\mathrm{own}-\mathrm{onl}+\mathrm{borr}$ holds because it credits borrowed shares.
\end{remark}

Scenario H shows where the collateral regime becomes visible: not in available inventory, but in profit and loss.

\begin{remark}[Pledge preserves economic ownership]
Under pledge (H), Bob retains economic ownership of posted collateral; the regime difference surfaces in PnL projection, where Bob's MSFT PnL is $c_p(200)\times P$, not $\mathrm{own}(0)\times P$.
\end{remark}

